\section{Introduction}
\label{sec:intro}

Let \(G,\Theta\in\mathbb R^{m\times n}\), with \(\Theta\ne0\).  Our primary
object is the nuclear-norm distance from \(G\) to the one-dimensional matrix
subspace generated by \(\Theta\):
\begin{equation}
 \phi(\lambda)=\|G+\lambda\Theta\|_*,\qquad
 d^\star=\min_{\lambda\in\mathbb R}\phi(\lambda)
 =\operatorname{dist}_{\|\cdot\|_*}
   \bigl(G,\operatorname{span}\{\Theta\}\bigr).
 \label{eq:D}
\end{equation}
Classical distance duality identifies the same value with
\begin{equation}
 p^\star=
 \max_{\Phi}
 \left\{\langle G,\Phi\rangle:
 \|\Phi\|_2\le1,\ \langle\Theta,\Phi\rangle=0\right\}.
 \label{eq:P}
\end{equation}
Thus \(d^\star=p^\star\): the dual unit ball is the spectral-norm ball, and
the annihilator of \(\operatorname{span}\{\Theta\}\) is one hyperplane.  We
study what happens when the affine matrix line
\(\lambda\mapsto G+\lambda\Theta\) meets the rank-deficient locus.  The
distinctive difficulty is that the natural point derivative is expressed by
the compact polar factor.  It is smooth on fixed-rank strata, but under
general rank-changing perturbations the zero-kernel completion need not have a
continuous extension; see, e.g., the classical polar-decomposition framework
of Higham~\cite{higham1986polar}.

When \(u,v\) are unit leading singular vectors at a simple top singular
value, \(\Theta=uv^\top\) gives the tangent linear maximization oracle
(LMO) on a smooth spectral level set.  Norm-ball LMOs underlie matrix-aware
first-order methods \citep{pethick2025lmo}.
Our motivating example is the Spectral Sphere Optimizer (SSO) of
Xie et al.~\cite[Section 3.2 and Appendix A.2, version 3]{xie2026controlled}, which uses a
scalar polar-factor equation; Li et al.~\cite[Appendix P.5]{li2026intrinsic} discuss the
same spectral-sphere subproblem.  The analysis below isolates the
rank-deficient certificate recovery needed for that equation.  It does not
assess the training results reported for SSO.

Related manifold oracles must be compared using their actual feasible sets.
The Stiefel product norm in Li et al.~\cite[Appendix H.3]{li2026intrinsic} and the
projected quotient norm of Yang and Lai~\cite{yang2026manifold} need not give the ambient
spectral ball intersected with the tangent space.  By contrast,
Solonko, Molozhavenko, and Rakhuba~\cite{solonko2026skewon} solve that joint-norm Stiefel problem exactly through
norm-preserving completion and explicitly include rank-deficient solution
families.  We make no general claim that tangent spectral LMOs lack closed
forms or that kernel completion itself is new.

\paragraph{Classical foundations and the scope of this work.}
Optimality in \eqref{eq:D} is classical Birkhoff--James orthogonality.
Distance duality, matrix-norm subdifferentials, and the associated criteria
are established in
\citep{singer1970best,watson1992subdifferential,zietak1993faces,
watson1993kyfan,bhatia1999orthogonality,grover2014distance,
johnston2022trace,grover2026kyfan}; we claim no novelty for those facts.
In particular, Liu and Vandenberghe~\cite{liu2009nuclear} already study nuclear-norm
minimization of general affine matrix-valued functions, its spectral-ball
dual, and Huber smoothing.  Hoheisel and Paquette~\cite{hoheisel2023uniqueness} characterize
flatness and uniqueness of affine nuclear-norm minimizers; uniqueness of
such an approximant is distinct from uniqueness of its constrained dual
certificate.  Our focus is the explicit classification and spectral centers
of the one-hyperplane slice
\(\partial\|A^\star\|_*\cap\Theta^\perp\), together with local
regularization-path formulas at its rank-changing points.

Nonsmooth singular-value calculus is classical
\citep{lewis2005nonsmooth1,lewis2005nonsmooth2}.
He, Kan, and Song~\cite{he2026twice} give second-order epi-derivatives of the nuclear norm,
including rank-deficient matrices.  These general variational results are
background for, rather than claims superseded by, our explicit
matrix-line centers, coefficients, cancellations, and uniform transition
laws.  Polar-factor sensitivity near rank loss is also established
\citep{higham1986polar,mathias1993polar}.

Canonical approximant selection by Schatten limits is also classical:
Legg and Ward~\cite{legg1985canonical} treat \(p\downarrow1\), while
Zi\k{e}tak~\cite{zietak2017strict} surveys the \(p\uparrow\infty\) strict-spectral
counterpart, including its conjectural general form.  Our selection instead
occurs on the equality-constrained certificate face.
If \(A=U_r\Sigma_rV_r^\top\) is a compact singular value decomposition,
the compact polar factor is
\begin{equation}
 \Sc(A)=U_rV_r^\top.
 \label{eq:compact-polar}
\end{equation}
At full rectangular rank, \(\phi\) is differentiable and
\(\phi'(\lambda)=\langle\Theta,\Sc(G+\lambda\Theta)\rangle\).  This suggests
the ordinary scalar equation
\begin{equation}
 h_{\rm c}(\lambda):=
 \langle\Theta,\Sc(G+\lambda\Theta)\rangle=0.
 \label{eq:compact-root}
\end{equation}
At full rectangular rank this equation is valid.  When the affine matrix loses
rank, it need not remain valid, because the complete subdifferential is
\begin{equation}
 \partial\|A\|_*
 =\left\{
 U_rV_r^\top+U_0KV_0^\top:\ \|K\|_2\le1
 \right\},
 \label{eq:full-subdiff}
\end{equation}
where \(U_0,V_0\) span the left and right null spaces.  The compact factor is
only the choice \(K=0\).  The omitted kernel block can be exactly what is
needed to enforce \(\langle\Theta,\Phi\rangle=0\).  Consequently,
\(h_{\rm c}\) can be monotone and still jump across zero.

Thus the LMO and nuclear-norm duality remain exact; only the replacement of a
set-valued optimality graph by one compact selection can fail.

\subsection{Contributions and novelty boundary}

The contributions, in the order used by the analysis, are as follows.

\begin{enumerate}
 \item Starting from the classical interval
 \(\partial\phi=[a-\beta,a+\beta]\)
 \citep{watson1992subdifferential,lewis1995convex}, we classify feasibility
 and uniqueness of the tangency slice.  Its minimum-Frobenius member follows
 by singular-value water-filling; clipping is universally unnecessary exactly
 for scalar multiples of partial isometries.

 \item An admissible Fenchel family unifies the hard and smooth centers:
 Huber gives \(K_F\), pseudo-Huber gives \(K_R\), and we characterize their
 agreement.  A polynomial derivative tail of order \(p>1\) yields exponent
 \(p/(p+1)\), while finite saturation gives a linear law.  The novelty beyond
 standard smoothing principles
 \citep{nesterov2005smooth,liu2009nuclear,lobos2025smooth,attouch1984variational} is the
 constrained spectral center and its coupling to the singular multiplier.

 \item We prove the regular expansion and, at every strict rank-deficient
 minimizer, analytic multiplier and returned-direction expansions without
 shape, normal-rank, or corank restrictions.  Compact-root instances have
 computable cubic/quadratic superconvergence and a quintic/quartic refinement.
 At transverse square corank-one boundaries, pseudo-Huber has the exact
 two-thirds law and a uniform cubic strict-to-boundary crossover.  Our new
 content is the matrix-line exponents, constants, and crossover, rather than
 the existence of fractional central-path rates \citep{basu2022central}.

 \item Finally, a strict \(2\times2\) witness extends to a nonempty open
 compact-polar failure set in every square dimension at least two, even with
 unique approximant and certificate.  This motivates retaining the full
 slice.  Rank-one normals also admit an SVD-free crossing formula.
\end{enumerate}

An open failure set implies neither an event-frequency law nor an outer
algorithmic or learning-performance gain; those require additional models.
For the rank-one SSO normal, the projected kernel normal has rank at most one,
so the hard and pseudo-Huber centers agree.  Their possible disagreement is
a feature of the more general matrix-line problem.

\subsection{Organization}

\Cref{sec:duality,sec:failure} give the scalar interval, structural failure,
and certificate geometry; \cref{sec:smoothing,sec:rates} develop the Fenchel
centers and local laws; \cref{sec:numerics} verifies the constants.  Auxiliary
level-set geometry and certificate repair are collected in the supplement.
